\documentclass[11pt,letterpaper]{article}

% Essential packages
\usepackage{fontspec}
\usepackage[utf8]{inputenc}
\usepackage[T1]{fontenc}
\usepackage{geometry}
\usepackage{amsmath}
\usepackage{amsfonts}
\usepackage{amssymb}
\usepackage{graphicx}
\usepackage{booktabs}
\usepackage{array}
\usepackage{longtable}
\usepackage{multirow}
\usepackage{multicol}
\usepackage{float}
\usepackage{caption}
\usepackage{subcaption}
\usepackage{hyperref}
\usepackage{fancyhdr}
\usepackage{datetime}
\usepackage{enumitem}
\usepackage{listings}
\usepackage{xcolor}
\usepackage{verbatim}
\usepackage{url}
\usepackage{svg}

% Page setup
\geometry{margin=1in}
\setlength{\parindent}{0pt}
\setlength{\parskip}{6pt}

% Header and footer setup
\pagestyle{fancy}
\fancyhf{}
\lhead{Data Science Capstone Journal}
\rhead{\today}
\cfoot{\thepage}
\renewcommand{\headrulewidth}{0.4pt}

% Code listing setup
\lstset{
    basicstyle=\ttfamily\footnotesize,
    backgroundcolor=\color{gray!10},
    frame=single,
    breaklines=true,
    captionpos=b,
    numbers=left,
    numberstyle=\tiny\color{gray},
    keywordstyle=\color{blue},
    commentstyle=\color{green!60!black},
    stringstyle=\color{red}
}

% Custom commands
\newcommand{\journalentry}[2]{
    \section*{Journal Entry - #1}
    \addcontentsline{toc}{section}{Journal Entry - #1}
    \textbf{Date:} #2 \\
    \textbf{Duration:} 10 hours \\
    \rule{\textwidth}{0.5pt}
}

\newcommand{\objective}[1]{
    \subsection*{Objective}
    #1
}

\newcommand{\activities}[1]{
    \subsection*{Activities Completed}
    #1
}

\newcommand{\findings}[1]{
    \subsection*{Key Findings}
    #1
}

\newcommand{\challenges}[1]{
    \subsection*{Challenges Encountered}
    #1
}

\newcommand{\nextSteps}[1]{
    \subsection*{Next Steps}
    #1
}

\newcommand{\reflections}[1]{
    \subsection*{Reflections}
    #1
}

% Title page setup
\newcommand{\workingtitle}{Statistical Truth Detection System}
\title{\Large \textbf{DATS 598: Data Science Capstone} \\ 
       \large Journal Entries and Progress Documentation \\
       \vspace{0.5cm}
       \normalsize \textbf{\workingtitle}}
\author{Eric Crisp \\ 
        University of Pennsylvania \\
        \texttt{ecrisp@upenn.edu}}
\date{\today}

\begin{document}

\maketitle
\tableofcontents
\newpage

% ========================================
% SAMPLE JOURNAL ENTRY -- MAIN PROPOSAL
% ========================================

\journalentry{Week 4 - \workingtitle}{\today}

\objective{
Complete initial data investigation, prototyping of architecture and product format, begin implementation.}

\activities{
\begin{itemize}
    \item Continued deep dive into the NLP processes needed, finding open source libraries where applicable, and determine the neccesary, nice to haves for project deliverables.
    \item 
\end{itemize}

}

\findings{



}

\nextSteps{

    \begin{itemize}
        \item Pull in data sources for fact-checking (e.g., FEVER, HuggingFace, Wikidata, etc.).
        
        \item Begin building out the backend applications via notebook (sentence parsing and analysis, API hooks, open source and python compability and configurations for environment stability).
        
        \item Begin prototyping the web app and setting up the tech stack (confirmation of CI/CD pipelines, containerized code, frontend-backend communication).
        
    \end{itemize}

    The current tech stack (as planned, and currently built is this):
    \begin{itemize}
        \item Backend: Python, FastAPI, Uvicorn
        \item Frontend: JavaScript, React, Vite
        \item VCS/Deployment: Git, GitHub Actions for CI/CD, AWS for hosting (?)
        \item Database: PostgreSQL, if RAG (?)
        \item NLP: SpaCy, NLTK, HuggingFace Transformers
        \item Containerization: Docker (Docker compose)
        \item EDA Packages: Jupyter Notebooks, Pandas, Matplotlib, Seaborn, NLTK, SpaCy, Transformers, Torch
    \end{itemize}
}

\end{document}